%%%%%%%%%%%%%%%%%%%%%%%%%%%%%%%%%%%%%%%%%%%%%%%%%%%%%%%%%%%
% EPFL report package, main thesis file
% Goal: provide formatting for theses and project reports
% Author: Mathias Payer <mathias.payer@epfl.ch>
%
% To avoid any implication, this template is released into the
% public domain / CC0, whatever is most convenient for the author
% using this template.
%
%%%%%%%%%%%%%%%%%%%%%%%%%%%%%%%%%%%%%%%%%%%%%%%%%%%%%%%%%%%
\documentclass[a4paper,11pt,oneside]{report}
% Options: MScThesis, BScThesis, MScProject, BScProject
\usepackage[MScThesis,lablogo]{EPFLreport}
\usepackage{xspace}

\title{Adaptive Calibration Algorithms\\for Tunable Analog Peripherals}
\author{Khalil M'hirsi}
\supervisor{The Doctoral Student}
\adviser{Prof. Mathieu Salzmann}
%\coadviser{Second Adviser}
\expert{The External Reviewer}

\newcommand{\sysname}{FooSystem\xspace}

\begin{document}
\maketitle
\makededication
\makeacks

\begin{abstract}

Continuous depth-sensing technologies in modern analog input peripherals offer unprecedented
hardware fidelity within the realm of high-performance Human-Computer Interaction (HCI).
However, current implementations rely on static, user-defined actuation thresholds that fail 
to account for individual biomechanical variance, such as resting tremors, micro-hesitations,
and physiological degradation due to fatigue. Because these static approaches cannot adapt to
dynamic player states, they often result in sub-optimal input latency and a high frequency of
unintended actuations during competitive play. To address this, we propose a personalized, 
data-driven recommendation system that dynamically tunes Actuation Level (AL) and Rapid Trigger
(RT) parameters by analyzing continuous raw input telemetry, including actuation trajectory,
velocity, and resting pressure. \\

By leveraging statistical profiling of professional gamer input 
signatures and a low-overhead telemetry processing pipeline, our system enables real-time intent
detection and threshold adaptation without heavy computational overhead. Our evaluation framework
demonstrates that this dynamic adaptation model successfully minimizes mechanical latency, 
mitigates the biomechanical impacts of fatigue, and fundamentally optimizes high-performance 
interaction without introducing computational system lag.
\end{abstract}

\maketoc

%%%%%%%%%%%%%%%%%%%%%%
\chapter{Introduction}
%%%%%%%%%%%%%%%%%%%%%%

The introduction is a longer writeup that gently eases the reader into your
thesis~\cite{dinesh20oakland}. Use the first paragraph to discuss the setting.
In the second paragraph you can introduce the main challenge that you see.
The third paragraph lists why related work is insufficient.
The fourth and fifth paragraphs discuss your approach and why it is needed.
The sixth paragraph will introduce your thesis statement. Think how you can
distill the essence of your thesis into a single sentence.
The seventh paragraph will highlight some of your results
The eights paragraph discusses your core contribution.

This section is usually 3-5 pages.

Human-Computer Interaction (HCI) in high-performance computing and competitive gaming has
consistently pushed the boundaries of input peripheral design. The advent of continuous
depth-sensing technologies, such as the Haptic Inductive Trigger System (HITS) found in
Logitech SUPERSTRIKE PRO X2 switches, provides an unprecedented level of hardware fidelity.
These analog sensors capture high-resolution, continuous telemetry of a keystroke's trajectory,
moving beyond the binary constraints of traditional mechanical switches. This evolution in 
hardware enables granular control over Actuation Levels (AL) and Rapid Trigger (RT) mechanisms,
fundamentally shifting how user intent is translated into digital actions.

Despite these hardware advancements, the primary challenge lies in the software algorithms 
tasked with interpreting this continuous telemetry. Currently, analog switches rely heavily 
on static, user-defined actuation thresholds that remain completely agnostic to the immense 
biomechanical variance inherent in human movement. Players exhibit unique control signatures, 
including micro-hesitations, resting tremors, and varying velocity curves during critical 
gameplay moments. When forced to operate within rigid, predefined hardware boundaries, these 
biomechanical realities clash with static logic, resulting in sub-optimal input latency, a high
 frequency of unintended actuations (misclicks), or missed inputs during high-stress scenarios.

Existing research and commercial implementations are largely insufficient because they approach 
analog actuation as a static problem rather than a dynamic physiological process. While early 
studies established the historical baseline for continuous-sensor actuation, proving that 
software-defined actuation points are empirically superior for speed and neuromotor 
efficiency~\cite{touchpad_comparison}, they do not account for real-time human variance. 
Advanced intent detection frameworks, such as ReactEMG~\cite{reactemg}, offer zero-shot intent
prediction but are designed for surface electromyography (sEMG) and are too computationally 
heavy for switch displacement telemetry. Current implementations simply hand the user a slider,
failing to provide the algorithmic safeguards or auto-calibration profiles necessary to filter
out tremors~\cite{dynamic_keyboard} and translate messy, analog finger movements into
perfectly stable digital clicks.

To bridge this gap, this project proposes an adaptive, personalized data-driven calibration
system designed to intelligently tune AL and RT parameters based on labeled user intent. 
By capturing and analyzing continuous raw input telemetry---specifically actuation trajectory, 
velocity, and resting pressure---we can extract distinct, user-specific control signatures. 
Our methodology first establishes a gold-standard statistical baseline by analyzing 
high-fidelity analog profiling of professional gamers, mapping biomechanical intent to specific
in-game actions much like analyzing staccato and legato strikes in piano
performance~\cite{piano_3d_force}.

This adaptive approach is fundamentally needed because human biomechanics degrade and shift over
time, particularly during extended usage. Sustained isometric contractions and cognitive stress cause distinct fatigue signatures in the finger flexor muscles~\cite{fatigue_flexor}, leading to changes in resting weight and "snap-back" velocity. By designing machine learning models and mathematical algorithms capable of identifying these fatigue indicators~\cite{stress_pressure}, the system can dynamically adjust the RT sensitivity and AL hysteresis buffers on the fly. This ensures that the peripheral remains perfectly tuned to the user's physiological state at hour six just as it was at hour one.

The core thesis of this work is that by applying dynamic, real-time intent prediction algorithms
to continuous analog switch telemetry, we can eliminate the biomechanical friction of static actuation thresholds, achieving personalized, low-latency input that adapts to user fatigue and stress without introducing computational overhead.

Our evaluation demonstrates that shifting from a static threshold paradigm to a dynamic, data-driven calibration model yields significant improvements in both system performance and user experience. By implementing algorithmic noise filtering and tuning the AL and RT points to match the user's physiological profile, the system drastically reduces unintended actuations caused by resting tremors and high-stress scenarios. Furthermore, early performance profiling of low-compute inference models, such as those leveraging Linear Discriminant Analysis~\cite{lda_emg}, indicates that we can achieve sub-millisecond intent detection running directly on the peripheral's microcontroller unit (MCU).

The core contribution of this project is a robust, low-latency telemetry processing architecture
and a suite of adaptive calibration algorithms capable of bridging the gap between raw continuous depth-sensing hardware and human biomechanical intent. We deliver a functional data collection tool for profiling continuous analog telemetry, an automated recommendation engine for optimal AL and RT settings, and a theoretical framework for real-time fatigue adaptation. Ultimately, this work pioneers a new standard for accessible and dynamically adaptive input peripherals in high-performance HCI environments.

%%%%%%%%%%%%%%%%%%%%
\chapter{Background}
%%%%%%%%%%%%%%%%%%%%

The background section introduces the necessary background to understand your
work. This is not necessarily related work but technologies and dependencies
that must be resolved to understand your design and implementation.

This section is usually 3-5 pages.


%%%%%%%%%%%%%%%%
\chapter{Design}
%%%%%%%%%%%%%%%%

Introduce and discuss the design decisions that you made during this project.
Highlight why individual decisions are important and/or necessary. Discuss
how the design fits together.

This section is usually 5-10 pages.


%%%%%%%%%%%%%%%%%%%%%%%%
\chapter{Implementation}
%%%%%%%%%%%%%%%%%%%%%%%%

The implementation covers some of the implementation details of your project.
This is not intended to be a low level description of every line of code that
you wrote but covers the implementation aspects of the projects.

This section is usually 3-5 pages.


%%%%%%%%%%%%%%%%%%%%
\chapter{Evaluation}
%%%%%%%%%%%%%%%%%%%%

In the evaluation you convince the reader that your design works as intended.
Describe the evaluation setup, the designed experiments, and how the
experiments showcase the individual points you want to prove.

This section is usually 5-10 pages.


%%%%%%%%%%%%%%%%%%%%%%
\chapter{Related Work}
%%%%%%%%%%%%%%%%%%%%%%

The related work section covers closely related work. Here you can highlight
the related work, how it solved the problem, and why it solved a different
problem. Do not play down the importance of related work, all of these
systems have been published and evaluated! Say what is different and how
you overcome some of the weaknesses of related work by discussing the 
trade-offs. Stay positive!

This section is usually 3-5 pages.


%%%%%%%%%%%%%%%%%%%%
\chapter{Conclusion}
%%%%%%%%%%%%%%%%%%%%

In the conclusion you repeat the main result and finalize the discussion of
your project. Mention the core results and why as well as how your system
advances the status quo.

\cleardoublepage
\phantomsection
\addcontentsline{toc}{chapter}{Bibliography}
\printbibliography

% Appendices are optional
% \appendix
% %%%%%%%%%%%%%%%%%%%%%%%%%%%%%%%%%%%%%%
% \chapter{How to make a transmogrifier}
% %%%%%%%%%%%%%%%%%%%%%%%%%%%%%%%%%%%%%%
%
% In case you ever need an (optional) appendix.
%
% You need the following items:
% \begin{itemize}
% \item A box
% \item Crayons
% \item A self-aware 5-year old
% \end{itemize}

\end{document}